The \peripheral{NPU} is a fixed-point neural network processing unit that accelerates the computation of a single fully-connected (dense) layer of a multilayer perceptron (MLP). The peripheral operates on data stored in a multiplexed 16\,KiB SRAM shared with the CPU. The CPU is responsible for loading input data and weights into the SRAM, configuring the NPU, and reading the results after the computation is complete.

\subsubsection*{Data Formats}

All NPU operands are fixed-point numbers, each stored in its own 32-bit SRAM word:

\begin{itemize}
    \item \textbf{Inputs}: signed Q0.15 format (16 bits: 1 sign bit + 15 fractional bits). Stored in bits 15:0 of a 32-bit SRAM word; bits 31:16 are unused.
    \item \textbf{Weights}: signed Q3.15 format (19 bits: 1 sign bit + 3 integer bits + 15 fractional bits). Stored in bits 18:0 of a 32-bit SRAM word; bits 31:19 are unused. Bit 18 is the sign bit.
    \item \textbf{Outputs}: signed Q3.15 format (19 bits: 1 sign bit + 3 integer bits + 15 fractional bits). Stored in bits 18:0 of a 32-bit SRAM word; bits 31:19 are unused. Bit 18 is the sign bit.
\end{itemize}

\subsubsection*{Memory Layout}

The input vector, weight matrix, and output vector must all reside within the same 16\,KiB multiplexed SRAM, and all three start addresses must be 4-byte aligned. Each operand (input, weight, or output) occupies exactly one 32-bit SRAM word. See the register descriptions for \texttt{NPUIVSAR}, \texttt{NPUWVSAR}, and \texttt{NPUOVSAR} for address details.

If bias is enabled (\texttt{NPUBEN}~=~1), the weight matrix contains \texttt{NPUNN}~+~1 rows of \texttt{NPUNI}~+~2 words each: one bias weight followed by \texttt{NPUNI}~+~1 synaptic weights per output neuron. If bias is disabled, each row contains \texttt{NPUNI}~+~1 synaptic weights.

\subsubsection*{Operation}

\begin{enumerate}
    \item Load the input vector into the SRAM at the desired address and write that address to \texttt{NPUIVSAR}.
    \item Load the weight matrix into the SRAM at the desired address and write that address to \texttt{NPUWVSAR}.
    \item Write the desired output address to \texttt{NPUOVSAR}.
    \item Configure \texttt{NPUCR}: set \texttt{NPUNI} to (number of inputs~$-$~1), \texttt{NPUNN} to (number of output neurons~$-$~1), and configure \texttt{NPUBEN} and \texttt{NPUAEN} as required.
    \item Set \texttt{NPUTHINK} to~1 to start the computation.
    \item Poll \texttt{NPUTHINK} until it reads~0, indicating that the computation is complete.
    \item Read the output vector from the SRAM at the address stored in \texttt{NPUOVA}.
\end{enumerate}
